\subsection{\texorpdfstring{关于 \(\lambda^{\#}\) 的另一种解释}{关于 \textbackslash lambda\^{}\{\textbackslash\#\} 的另一种解释}}

对于每个 \(x \in X\)，从 H 到 X 的映射 \(\xi \mapsto x\xi\) 是 逆紧的（GT, III, §4, No.~2, Prop. 4），因此 \(\beta\) 在此映射下在 X 上有一个像测度，该像测度集中于轨道 xH（Ch. V, §6, No.~2, Cor. 3 of Prop. 2）；由于 \(\beta\) 左不变，这个像测度只依赖于 x 在 X/H 中的类 \(u = \pi(x)\)，记作 \(\beta_{u}\)。按定义，对于 \(f \in \mathcal{K}(X)\)，

\[
\int_ {\mathrm{X}} f (y) d \beta_ {u} (y) = \int_ {\mathrm{H}} f (x \xi) d \beta (\xi) = f ^ {\flat} (u).\tag{8}
\]

由此可见

\[
(\varepsilon_ {u}) ^ {\sharp} = \beta_ {u}.\tag{9}
\]

引理 2. --- 设 \(f\) 是 \(X\) 上的函数，取值于一个拓扑空间。

\begin{enumerate}
\def\labelenumi{\alph{enumi})}
\tightlist
\item
  若 \(f\) 是一个 \(\geqslant 0\) 的数值函数，则对于 \(x \in X\)，
\end{enumerate}

\[
\int_ {X} ^ {*} f (y) d \beta_ {\dot {x}} (y) = \int_ {H} ^ {*} f (x \xi) d \beta (\xi) \quad (\dot {x} = \pi (x)).
\]

\begin{enumerate}
\def\labelenumi{\alph{enumi})}
\setcounter{enumi}{1}
\item
  \(f\) 为 \(\beta_{\dot{x}}\)-可测的充要条件是，函数 \(\xi \mapsto f(x\xi)\) 在 \(H\) 上为 \(\beta\)-可测的。
\item
  设 \(\mathbf{f}\) 是 \(\mathbf{X}\) 上的函数，取值于 Banach 空间或 \(\overline{\mathbf{R}}\)；则，\(\mathbf{f}\) 为 \(\beta_{\dot{x}}\)-可积（分别为本质上 \(\beta_{\dot{x}}\)-可积）的充要条件是，函数 \(\xi \mapsto \mathbf{f}(x\xi)\) 在 \(\mathbf{H}\) 上为 \(\beta\)-可积（分别为本质上 \(\beta\)-可积）；在此情形下，\(\int_{\mathbf{X}} \mathbf{f}(y) d\beta_{\dot{x}}(y) = \int_{\mathbf{H}} \mathbf{f}(x\xi) d\beta(\xi)\)。
\end{enumerate}

这由 Ch. V, §4, Prop. 2、Prop. 3 和 Th. 2 得出。

由于 \(f^{\flat} \in \mathcal{K}(X/H)\) 对 \(f \in \mathcal{K}(X)\) 成立，公式 (8) 证明了：映射 \(u \mapsto \beta_{u}\) 从 X/H 到 \(\mathcal{M}(X)\) 是弱连续的；族 \((\beta_{u})\) 是 \(\lambda\)-适合的 \(^{1}\)，对于 X/H 上任意正测度 \(\lambda\)；并且

\[
\lambda^ {\sharp} = \int_ {\mathrm{X/H}} \beta_ {u} d \lambda (u),\tag{10}
\]

这给出了 \(\lambda^{\sharp}\) 的一种新解释。

命题 5. --- 设 \(\lambda\) 是 X/H 上的正测度。

\begin{enumerate}
\def\labelenumi{\alph{enumi})}
\item
  设 \(f\) 是一个 \(\lambda^{\sharp}\)-可测函数，定义在 \(X\) 上，取值于一个拓扑空间，并且在某个 \(\lambda^{\sharp}\)-可积集的可数并之外为常值。则那些 \(\dot{x} \in X / H\) 使函数 \(\xi \mapsto f(x\xi)\) 不是 \(\beta\)-可测的点所成的集合是局部 \(\lambda\)-可忽略的。
\item
  设 \(f\) 是一个 \(\lambda^{\sharp}\)-可测的 \(\geqslant 0\) 函数，定义在 \(X\) 上，并在某个 \(\lambda^{\sharp}\)-可积集的可数并之外为零。则函数 \(\dot{x} \mapsto \int^{*} f(x\xi) d\beta(\xi)\) 在 \(X / H\) 上是 \(\lambda\)-可测的，并且
\end{enumerate}

\[
\int_ {\mathrm{X}} ^ {*} f (x) d \lambda^ {\sharp} (x) = \int_ {\mathrm{X/H}} ^ {*} d \lambda (\dot {x}) \int_ {\mathrm{H}} ^ {*} f (x \xi) d \beta (\xi) \qquad \big (\dot {x} = \pi (x) \big).
\]

\begin{enumerate}
\def\labelenumi{\alph{enumi})}
\setcounter{enumi}{2}
\tightlist
\item
  设 \(\mathbf{f}\) 是一个 \(\lambda^{\sharp}\)-可积函数，定义在 \(X\) 上，取值于 Banach 空间或 \(\overline{\mathbf{R}}\)。则那些 \(\dot{x} \in X / H\) 使得函数 \(\xi \mapsto \mathbf{f}(x\xi)\) 不是 \(\beta\)-可积的点所成的集合是 \(\lambda\)-可忽略的；由下式几乎处处定义的函数 \(\mathbf{f}^{\flat}\) 在 \(X / H\) 上
\end{enumerate}

\[
\mathbf {f} ^ {\flat} (\dot {x}) = \int_ {\mathrm{H}} \mathbf {f} (x \xi) d \beta (\xi) \quad \left(\dot {x} = \pi (x)\right)\tag{11}
\]

是 \(\lambda\)-可积的，并且

\[
\int_ {\mathrm{X/H}} \mathbf {f} ^ {\flat} d \lambda = \int_ {\mathrm{X}} \mathbf {f} d \lambda^ {\sharp}\tag{12}
\]

以及

\[
\int_ {\mathrm{X} / \mathrm{H}} | \mathbf {f} ^ {\flat} | d \lambda \leqslant \int_ {\mathrm{X}} | \mathbf {f} | d \lambda^ {\sharp}.\tag{13}
\]

\begin{enumerate}
\def\labelenumi{\alph{enumi})}
\setcounter{enumi}{3}
\tightlist
\item
  设 \(\mathbf{f}\) 是一个 \(\lambda^{\sharp}\)-可测函数，定义在 \(\mathbf{X}\) 上，取值于 \(\overline{\mathbf{R}}\) 或 Banach 空间，并且在某个 \(\lambda^{\sharp}\)-可积集的可数并之外为零。则 \(\mathbf{f}\) 为 \(\lambda^{\sharp}\)-可积的充要条件是
\end{enumerate}

\[
\int_ {\mathrm{X} / \mathrm{H}} ^ {*} d \lambda (\dot {x}) \int_ {\mathrm{H}} ^ {*} | \mathbf {f} (x \xi) | d \beta (\xi) <   + \infty \quad \left(\dot {x} = \pi (x)\right).
\]

结合引理 2，断言 a)、b) 和 c) 由 Ch. V, §3, Prop. 4、Prop. 5 和 Th. 1 得出（(13) 除外；(13) 由 (12) 得出，因为显然有 \(|\mathbf{f}^{\flat}| \leqslant |\mathbf{f}|^{\flat}\)）；d) 由 b) 得出。

命题 6. --- 设 \(\lambda\) 是 X/H 上的正测度。

\begin{enumerate}
\def\labelenumi{\alph{enumi})}
\item
  设 N 是 X/H 的子集。N 局部 \(\lambda\)-可忽略的充要条件是，\(\overline{\pi}^{-1}(N)\) 局部 \(\lambda^{\sharp}\)-可忽略。
\item
  设 \(g\) 是 \(\mathrm{X / H}\) 上取值于一个拓扑空间的函数。\(g\) 为 \(\lambda\)-可测的充要条件是 \(g \circ \pi\) 为 \(\lambda^{\sharp}\)-可测的。
\item
  设 \(\mathbf{h}\) 是 \(\mathrm{X / H}\) 上取值于 Banach 空间或 \(\overline{\mathbf{R}}\) 的函数。\(\mathbf{h}\) 局部 \(\lambda\)-可积的充要条件是 \(\mathbf{h} \circ \pi\) 局部 \(\lambda^{\sharp}\)-可积；在此情形下，\((\mathbf{h} \cdot \lambda)^{\sharp} = (\mathbf{h} \circ \pi) \cdot \lambda^{\sharp}\)。
\end{enumerate}

假设 \(\mathbf{h} \circ \pi\) 局部 \(\lambda^{\sharp}\)-可积。对于每个 \(f \in \mathcal{K}(X)\)，\(f \cdot (\mathbf{h} \circ \pi)\) 都是 \(\lambda^{\sharp}\)-可积的，因此由 Prop. 5，函数 \((f \cdot (\mathbf{h} \circ \pi))^b = f^b \cdot \mathbf{h}\) 是 \(\lambda\)-可积的，并且

\[
\int_ {\mathrm{X} / \mathrm{H}} f ^ {\flat} \cdot \mathbf {h} d \lambda = \int_ {\mathrm{X}} f \cdot (\mathbf {h} \circ \pi) d \lambda^ {\sharp}.
\]

由于 \(f \mapsto f^{\flat}\) 是从 \(\mathcal{K}(\mathrm{X})\) 到 \(\mathcal{K}(\mathrm{X}/\mathrm{H})\) 上的满射，这表明 \(\mathbf{h}\) 局部 \(\lambda\)-可积，并且

\[
(\mathbf {h} \cdot \lambda) ^ {\sharp} = (\mathbf {h} \circ \pi) \cdot \lambda^ {\sharp}.
\]

特别地，若 \(\overline{\pi}^{1}(N)\) 局部 \(\lambda^{\sharp}\)-可忽略，则 \(\varphi_{N} \circ \pi\) 局部 \(\lambda^{\sharp}\)-可忽略，因此 \((\varphi_{N} \cdot \lambda)^{\sharp} = (\varphi_{N} \circ \pi) \cdot \lambda^{\sharp} = 0\)，从而 \(\varphi_{N} \cdot \lambda = 0\)，且 \(N\) 局部 \(\lambda\)-可忽略。现在假设 \(g \circ \pi\) 是 \(\lambda^{\sharp}\)-可测的。设 \(K'\) 是 \(X/H\) 的紧子集。取 \(f \in \mathcal{K}_{+}(X)\)，使得 \(f^{\flat} = 1\) 在 \(K'\) 上（No.~1, Prop. 2），并令 \(K = \operatorname{Supp} f\)；有 \(\pi(K) \supset K'\)。存在 \(K\) 的一个分割，它由一个 \(\lambda^{\sharp}\)-可忽略集 \(M\) 和紧集序列 \((K_n)\) 构成，使得 \((g \circ \pi)|K_n\) 连续，对每个 \(n\) 成立。于是 \(g|\pi(K_n)\) 连续。令 \(P\) 为 \(K'\) 中不属于 \(\pi(K_1) \cup \pi(K_2) \cup \ldots\) 的点集；则 \(\overline{\pi}^{1}(P) \cap K\) 包含于 \(M\) 中，因而是 \(\lambda^{\sharp}\)-可忽略的；所以 \(f \cdot \varphi_{\overline{\pi}(P)}^{-1}\) 是 \(\lambda^{\sharp}\)-可忽略的；由此（Prop. 5）可得

\[
0 = \int_ {\mathrm{X}} f \cdot \varphi_ {\pi^ {- 1} (\mathrm{P})} d \lambda^ {\sharp} = \int_ {\mathrm{X/H}} f ^ {\flat} \cdot \varphi_ {\mathrm{P}} d \lambda \geqslant \int_ {\mathrm{X/H}} ^ {*} \varphi_ {\mathrm{P}} d \lambda ,
\]

因此 \(\mathbf{P}\) 是 \(\lambda\)-可忽略的，且 \(g\) 是 \(\lambda\)-可测的。

若 N 局部 \(\lambda\)-可忽略，则 \(\overline{\pi}^{1}(N)\) 局部 \(\lambda^{\sharp}\)-可忽略（Ch. V, §6, No.~6, Cor. 1 of Prop. 10）。若 g 是 \(\lambda\)-可测的，则 \(g \circ \pi\) 是 \(\lambda^{\sharp}\)-可测的（同上）。最后，假设 h 局部 \(\lambda\)-可积。我们已经知道 \(h \circ \pi\) 是 \(\lambda^{\sharp}\)-可测的。对于每个 \(f \in \mathcal{K}_{+}(X)\)，由 Prop. 5 有

\[
\int_ {X} ^ {*} f (x) | \mathbf {h} | (\pi (x)) d \lambda^ {\sharp} (x) = \int_ {X / H} ^ {*} | \mathbf {h} | (u) f ^ {\flat} (u) d \lambda (u) <   + \infty ,
\]

因此 \(\mathbf{h} \circ \pi\) 局部 \(\lambda^{\sharp}\)-可积。

推论 1. --- 设 \(\lambda, \lambda'\) 是 X/H 上的两个正测度。\(\lambda'\) 以 \(\lambda\) 为基的充要条件是 \(\lambda'^{\sharp}\) 以 \(\lambda^{\sharp}\) 为基。\(\lambda\) 与 \(\lambda'\) 等价的充要条件是 \(\lambda^{\sharp}\) 与 \(\lambda'^{\sharp}\) 等价。

第一个断言由 Prop. 6 的 a) 和 c) 得出。第二个断言由第一个断言得出。

推论 2. --- 设 \(\lambda\) 是 X/H 上的正测度，并设 \(f\) 是 X 上的 \(\lambda^{\sharp}\)-可测数值函数。假设对于每个 \(\xi \in H\)，都有 \(\delta(\xi)f = f\) 局部 \(\lambda^{\sharp}\)-几乎处处成立。则存在 X/H 上的 \(\lambda\)-可测函数 \(g\)，使得 \(f = g \circ \pi\) 局部 \(\lambda^{\sharp}\)-几乎处处成立。

将 \(f\) 替换为 \(f / (1 + |f|)\) 后，可归结为 \(f\) 有界的情形，因而局部 \(\lambda^{\sharp}\)-可积。令 \(\mu = f \cdot \lambda^{\sharp}\)。关于 \(f\) 的假设蕴含，有 \(\pmb{\delta}(\xi) \mu = f \cdot \pmb{\delta}(\xi) \lambda^{\sharp} = \Delta_{\mathrm{H}}(\xi) \mu\)，对于所有 \(\xi \in \mathrm{H}\) 成立。于是存在（Prop. 4）\(\lambda'\)，它是 \(\mathrm{X} / \mathrm{H}\) 上的测度，使得 \(\mu = \lambda'^{\sharp}\)。由推论 1，关于 \(\lambda\)，存在函数 \(g\)，它定义在 \(\mathrm{X} / \mathrm{H}\) 上且局部可积，使得 \(\lambda' = g \cdot \lambda\)。由 Prop. 6，\(f \cdot \lambda^{\sharp} = \lambda'^{\sharp} = (g \circ \pi) \cdot \lambda^{\sharp}\)，从而 \(f = g \circ \pi\) 局部 \(\lambda^{\sharp}\)-几乎处处成立。

推论 3. --- a) 设 \((\lambda_{\iota})_{\iota \in I}\) 是 X/H 上的实测度族。族 \((\lambda_{\iota})\) 在 \(\mathcal{M}(\mathrm{X / H})\) 中有上界的充要条件是族 \((\lambda_{\iota}^{\sharp})\) 在 \(\mathcal{M}(\mathrm{X})\) 中有上界；在此情形下

\[
\sup (\lambda_ {\iota} ^ {\sharp}) = (\sup \lambda_ {\iota}) ^ {\sharp}.
\]

\begin{enumerate}
\def\labelenumi{\alph{enumi})}
\setcounter{enumi}{1}
\item
  设 \(\lambda\) 是 X/H 上的实测度。则 \((\lambda^{+})^{\sharp} = (\lambda^{\sharp})^{+}\)，且 \((\lambda^{-})^{\sharp} = (\lambda^{\sharp})^{-}\)。
\item
  设 \(\lambda\) 是 X/H 上的复测度。则 \(|\lambda|^{\sharp} = |\lambda^{\sharp}|\)。
\end{enumerate}

假设族 \((\lambda_{\iota})\) 有上界，并令 \(\mu = \sup \lambda_{\iota}\)。由于 \(\lambda \geqslant 0\) 蕴含 \(\lambda^{\sharp} \geqslant 0\)，故有 \(\mu^{\sharp} \geqslant \lambda_{\iota}^{\sharp}\)，对所有 \(\iota\) 成立，这表明族 \((\lambda_{\iota}^{\sharp})\) 有上界，并且

\[
(\sup \lambda_ {\iota}) ^ {\sharp} \geqslant \sup (\lambda_ {\iota} ^ {\sharp}).
\]

反之，假设族 \((\lambda_{\iota}^{\sharp})\) 有上界，并令 \(\nu = \sup(\lambda_{\iota}^{\sharp})\)。由于有 \(\delta(\xi)\lambda_{\iota}^{\sharp} = \Delta_{\mathrm{H}}(\xi)\lambda_{\iota}^{\sharp}\)，对于所有 \(\xi \in H\) 成立，显然 \(\delta(\xi)\nu = \Delta_{\mathrm{H}}(\xi)\nu\)，因此存在 \(\mu' \in \mathcal{M}(\mathrm{X}/\mathrm{H})\)，使得 \(\nu = \mu'^{\sharp}\)。由于 \(\lambda^{\sharp} \geqslant 0\) 蕴含 \(\lambda \geqslant 0\)，故 \(\mu' \geqslant \lambda_{\iota}\)，对所有 \(\iota\) 成立，这表明族 \((\lambda_{\iota})\) 有上界，并且 \(\nu = \mu'^{\sharp} \geqslant (\sup \lambda_{\iota})^{\sharp}\)，从而

\[
\sup \left(\lambda_ {\iota} ^ {\sharp}\right) \geqslant \left(\sup \lambda_ {\iota}\right) ^ {\sharp},
\]

这就完成了 a) 的证明。b) 的断言立即得出，因为例如 \(\lambda^{+}\) 正是 \(\sup(\lambda,0)\)。为证明 c)，只需注意，\(|\lambda| = \sup \mathcal{R}(\alpha \lambda)\)，其中上确界取遍绝对值为 1 的复数 \(\alpha\)；另一方面，\(\mathcal{R}(\mu^{\sharp}) = (\mathcal{R}\mu)^{\sharp}\) 对于每个 \(\mu \in \mathcal{M}(\mathrm{X / H})\) 成立。

注记。--- 1) Prop. 6 a) 可表述为：\(\lambda\) 是 \(\lambda^{\sharp}\) 在 \(\pi\) 下的拟像测度（Ch. VI, §3, No.~2, Def. 1）。

\begin{enumerate}
\def\labelenumi{\arabic{enumi})}
\setcounter{enumi}{1}
\tightlist
\item
  假设 H 紧且 \(\beta\) 已归一化。X 的每个紧子集的饱和集都是紧的。因此，若 \(f \in \mathcal{K}(\mathrm{X}/\mathrm{H})\)，则 \(f \circ \pi \in \mathcal{K}(\mathrm{X})\)；并且对于 X/H 上的每个正测度 \(\lambda\)，Prop. 5 c) 给出
\end{enumerate}

\[
\int_ {\mathrm{X}} (f \circ \pi) (x) d \lambda^ {\sharp} (x) = \int_ {\mathrm{X/H}} f (u) d \lambda (u).
\]

换言之，\(\lambda\) 是 \(\lambda^{\sharp}\) 在 \(\pi\) 下的像测度。

\begin{enumerate}
\def\labelenumi{\arabic{enumi})}
\setcounter{enumi}{2}
\item
  Prop. 6 的 Cor. 3 c) 立即表明，本小节的结果在复测度情形仍然成立（涉及上积分的结果除外）。
\item
  设 \(\mathbf{m}\) 是 X/H 上取值于 E 的向量测度，且 \(q\) 是 E 上的下半连续半范数。\(\mathbf{m}\) 为 \(q\)-可控的（Ch. VI, §2, No.~3, Def. 3）的充要条件是 \(\mathbf{m}^{\sharp}\) 具有同样性质；在此情形下 \(q(\mathbf{m}^{\sharp}) = q(\mathbf{m})^{\sharp}\)。这立即由定义和 Cor. 3 a) 得出。
\end{enumerate}

另一方面，设 \(\mu\) 是 X/H 上的正测度。m 按标量以 \(\mu\) 为基的充要条件是 \(\mathbf{m}^{\sharp}\) 按标量以 \(\mu^{\sharp}\) 为基；这由 Cor. 1 得出。

最后，若 \(\mathbf{m}\) 以 \(\mu\) 为基，并且密度为 \(\mathbf{f}\)（关于 \(\mu\)，Ch. VI, §2, No.~4, Def. 4），则 \(\mathbf{m}^{\sharp}\) 以 \(\mu^{\sharp}\) 为基，密度为 \(\mathbf{f} \circ \pi\)：这由 Prop. 6 c) 得出。

